% Discussion + Limitations. Discussion subsection added v2: the
% basin-selection mechanism revision is the central interpretive claim
% of the paper and previously lived only as a paragraph inside Phase 2C.
% Pulled forward here so the abstract and introduction can reference a
% dedicated home for it.

\section{Discussion and Limitations}
\label{sec:limitations}

\subsection{Discussion}
\label{sec:discussion}

\paragraph{Adaptive O2O wins via basin selection, not adaptive trading.}
The pre-registered hypothesis behind the adaptive conservatism schedule
(Sec.~\ref{sec:method:o2o}) was that the modulated $\beta_t$ would let
the SAC fine-tune respond to distribution shift in the test window:
relax pessimism when the online state distribution looks in-support,
saturate it when not. The Phase~2C results disconfirm that mechanism.
Both naive and adaptive O2O converge to functionally static allocations
(per-seed turnover $\le 2.2\!\times\!10^{-4}$ across all six runs ---
three orders of magnitude below GRPO online's $0.24$ and two below
behavior-anchored offline RL's $\sim 0.002$); neither method is doing
anything resembling regime-aware rebalancing on the 1{,}061-day test
window. The differential between the two is therefore not how they
trade --- it is which static allocation the SAC fine-tune lands on.
Naive O2O's pinned $\beta_t \equiv \alpha_{\mathrm{CQL}}$ leaves the
basin random across seeds (Sharpe range $[-0.236,\,+2.179]$, std $1.21$;
$0.03$rd to $100$th percentile of the static-allocation reference
distribution in Fig.~\ref{fig:static-sharpe-cdf} --- effectively a
random sample from the full simplex); adaptive O2O's modulated schedule
consistently steers the same SAC dynamics from the same Geodesic-CQL
pretrain to a uniformly strong basin (Sharpe range
$[+1.629,\,+1.931]$, std $0.15$, percentile band
$[98.74, 99.93]$): an eight-fold reduction in seed-to-seed Sharpe std
and a sharper narrowing in percentile space (full simplex
distribution $\rightarrow$ top $1.25\%$). The conservatism schedule's load-bearing
function in our results is convergence-basin steering during fine-tune,
not online responsiveness during deployment. We surface this as a
mechanism revision rather than a confirmation, in keeping with the
post-hoc-vs-pre-registered distinction; the empirical headline (adaptive
beats classical baselines, n=3 seeds, std $0.15$) is unchanged.

\paragraph{Implications for O2O scheduling research.}
The standard framing of O2O conservatism scheduling research is
``control online responsiveness to distribution shift'' --- the
schedule's job is to govern how aggressively the online phase departs
from the offline prior as the deployment distribution drifts. Our
finding suggests a complementary axis on which the same schedule is
load-bearing: it shapes which fixed point of the SAC fine-tune dynamics
the policy converges to, with little observable activity at deployment
time. For domains where SAC fine-tune dynamics are well-identified
(rich online reward signal, small batch noise, unconstrained continuous
actions), the two framings collapse onto the same behavior. For domains
where SAC fine-tune is weakly identified --- small batches, low
signal-to-noise, simplex-bounded actions, regimes where the
auto-temperature mechanism degenerates --- basin steering may be the
more useful design target for the schedule, and the relevant evaluation
is seed-stability of the converged policy rather than online tracking
of distribution drift. We do not have evidence that this generalizes
beyond the 8-ETF allocation setting; we report it as a hypothesis
generated by the data, to be tested in future O2O work where the
SAC fine-tune is similarly weakly identified.

\paragraph{Continuous GRPO: sound formulation, dominated equilibrium on
this universe.} The Continuous GRPO formulation (Sec.~\ref{sec:method:grpo})
is methodologically sound on this task: training is stable, the
group-size ablation reveals a clean compute-vs-precision tradeoff
(mean Sharpe and inter-seed std both decrease with $G$;
Table~\ref{tab:phase2d}), the critic-free advantage stays finite at
constant-reward states by construction, and the exogeneity precondition
is unit-tested at the env boundary. The formulation does what we ask
of it. The empirical answer it produces, however, is dominated. GRPO
online converges to an active-trading equilibrium with turnover $\sim
0.24$, which on this 8-ETF universe at $\lambda{=}0.001$ pays a per-step
friction tax that no daily directional signal recovers; the result is
test Sharpe $0.297$, well below the equal-weight $0.953$ reference. We
read this as a property of the universe, not of the algorithm: at this
asset count, frequency, and friction level, the static equal-weight
allocation is a strong dominator that any active-trading method must
outperform after costs, and our empirical evidence is that no online
method we ran clears that bar. We frame Continuous GRPO as a
methodological contribution --- a critic-free formulation that exploits
exogeneity --- and explicitly do not claim it beats classical baselines
on retail allocation under linear cost. The two claims are independent;
domains with non-exogenous transitions, sparser reward, or weaker
classical dominators are where the formulation might be evaluated
empirically.

\subsection{Limitations}
\label{sec:limit}

\paragraph{Single test window.} The chronological-split protocol gives
us exactly one test window: 2022-Q1 through 2026-Q1. Walk-forward
validation across multiple test windows would be the right measurement
of method robustness; we run a single window because the milestone's
2022 stock-bond correlation break is the regime change the experimental
design was built around. With one window, every reported Sharpe number
is a single sample from an unknown sampling distribution. We report
seed-variance because that is the variance the experiment can produce;
we do not claim that the per-method ranking generalizes to unseen test
windows of the same length.

\paragraph{Behavior-mixture sensitivity.} The offline buffer is rolled
under a 4-way mixture of Dirichlet, equal-weight, momentum, and
risk-parity policies (Sec.~\ref{sec:experiments:setup}). Earlier
results in this codebase (the milestone's 27-IQL ablation) used a
uniform-Dirichlet buffer instead, and we discovered during Phase 2 prep
that the momentum component had been numerically degenerate
(Appendix~\ref{app:implementation}). Both the
behavior policy and the IQL $\lambda$-anchor in Phase 2A serve to
quantify how much the new mixture changes the apples-to-apples
comparison; we do not claim invariance to behavior-policy choice. A
follow-up systematically varying the mixture composition is the natural
next step.

\paragraph{Transaction-cost model.} Eq.~\ref{eq:reward} models linear
turnover cost only --- no slippage, market impact, or asymmetric
bid/ask. Realistic execution friction at the 8-ETF retail-portfolio
scale is plausibly \emph{lower} than $\lambda{=}0.001$ for SPY-like
names and \emph{higher} for thinly-traded ones (UUP, HYG); a per-asset
$\lambda$ vector would be a strict refinement. We report at three
flat-$\lambda$ operating points to bracket the realistic range and flag
the modeling choice rather than hide behind a single ``representative''
number.

\paragraph{Seed reproducibility.} The PortfolioEnv random-start RNG is
constructed at \texttt{core/envs/portfolio\_env.py:87} without reading
the global numpy seed. Torch model init, action sampling, and
replay-buffer sampling are correctly seeded; env-side episode-start
randomness is not. Reported seed-variance is therefore ``per-process
variance, env-randomness uncontrolled'' rather than bit-exact
seed-replicable variance. The fix is one-line; we do not apply it
mid-run because doing so would invalidate the in-flight metrics.json
results. Post-paper follow-up.

\paragraph{Continuous-GRPO scope.} Our critic-free formulation
(Sec.~\ref{sec:method:grpo}) is sound iff the env satisfies the
exogeneity property in Eq.~\ref{eq:exogeneity}. We provide unit tests
that make the precondition legible
(\texttt{tests/test\_exogeneity.py}); we do not claim the method
transfers to settings where actions affect future state. In particular,
trades large enough to move the underlying market, multi-period option
strategies, and any setting with persistent inventory state would
violate the precondition. Any port of Continuous GRPO to a new env
should ship with its own analogue of the exogeneity test.

\paragraph{Adaptive-O2O bug fix.} The pre-Phase-2 codebase silently
ran the ``adaptive O2O'' code path with conservatism turned off
(Appendix~\ref{app:implementation}). Every adaptive-vs-naive number we
report post-dates the fix. Comparison with any prior milestone-attributed
adaptive-O2O number is therefore not apples-to-apples; we surface this
here rather than rewrite the milestone narrative.

\paragraph{Post-paper follow-up: theoretical extensions on critic
removal under exogeneity.} The exogeneity argument that retires the
value baseline (Sec.~\ref{sec:method:grpo:exogeneity}) suggests two
formal extensions: (a) a fully critic-free actor-only policy gradient
under action-independent dynamics, with no Bellman backup, no TD
learning, and analytically tractable variance properties; (b) a
Rao-Blackwellized estimator that marginalizes the action-independent
future randomness, with provably lower variance than REINFORCE under
the exogeneity precondition. Both are theoretical contributions out of
scope for this paper's empirical focus, and both are constrained to
settings where the precondition holds.
